% Method description; no performance values are embedded.
Both variants use the same 96-dimensional, three-layer patch Transformer.
\emph{CorA-MAPLE-E} samples historical excerpts, whereas \emph{CorA-MAPLE-F}
uses complete unlabelled flood episodes. The student learns from masked and
corrupted inputs while the moving-average teacher processes the unchanged
reference. Pretraining combines latent prediction, balanced binary reconstruction,
and variance regularization, weighted by $1$, $0.50$, and $0.10$, respectively,
with teacher momentum $0.995$.

After pretraining, the student encoder and a new linear classification head are
jointly fine-tuned using only clean labelled observations:
\begin{equation}
L_{\mathrm{cls}}=-\frac{1}{N_L}\sum_{i=1}^{N_L}\log p_{y_i}(X_i),
\label{eq:clean_supervised}
\end{equation}
where $N_L$ is the number of labelled episodes, $y_i$ is the reference class, and
$p_{y_i}(X_i)$ is its predicted probability. No experimental corruption or patch
masking is applied during supervised training.

Full-channel false activation, full-channel loss, and timestamp uncertainty are
sampled only during corruption-aware pretraining. Delayed observations and missing
or additional ACT/RTN notifications are reserved for testing. Event reconstruction
uses the recorded per-channel initial state as an independent snapshot. Delayed
windows retain explicit valid lengths, excluding padded samples from attention,
pooling and the classical reference representations.
